% Largest per-axis pointing jitter at which a 290 GHz or a 1550 nm crosslink holds a target rate with
% 99.9 percent availability and 3 dB of margin (0.3 m apertures, 1 W, LCRNS mean and worst ranges).
% Data: scripts/jitter.py, from the link budget of the space-segment companion. Drawn for the tutorial.
\begin{tikzpicture}
\begin{axis}[
  width=0.62\textwidth, height=5.4cm,
  xmode=log, ymode=log, xmin=1, xmax=5e4, ymin=0.5, ymax=3000,
  xlabel={Target information rate (Mbit/s)}, ylabel={Tolerable per-axis jitter $\sigma$ (\textmu rad)},
  grid=major, grid style={gray!20}, tick label style={font=\scriptsize}, label style={font=\footnotesize},
  unbounded coords=jump,
  legend style={font=\scriptsize, at={(1.03,0.5)}, anchor=west, draw=gray!50, cells={anchor=west}, row sep=1pt},
]
\addplot[lowc, very thick] table[x=R, y=rf_mean] {figs/jitter_tol.dat};
\addlegendentry{\SI{290}{GHz}, mean range}
\addplot[lowc, thick, densely dashed] table[x=R, y=rf_worst] {figs/jitter_tol.dat};
\addlegendentry{\SI{290}{GHz}, worst range}
\addplot[highc, very thick] table[x=R, y=op_mean] {figs/jitter_tol.dat};
\addlegendentry{\SI{1550}{nm}, mean range}
\addplot[highc, thick, densely dashed] table[x=R, y=op_worst] {figs/jitter_tol.dat};
\addlegendentry{\SI{1550}{nm}, worst range}
\draw[black!45, dash dot] (axis cs:1,50) -- (axis cs:5e4,50) node[pos=0.98, anchor=south east, font=\scriptsize, text=black!60] {star tracker, \SI{50}{\micro\radian}};
\draw[black!45, dash dot] (axis cs:1,1) -- (axis cs:5e4,1) node[pos=0.02, anchor=north west, font=\scriptsize, text=black!60] {precision pointing, \SI{1}{\micro\radian}};
\fill[black] (axis cs:62,2.65) circle (1.6pt);
\node[anchor=north west, font=\scriptsize] at (axis cs:70,2.5) {crossover};
\end{axis}
\end{tikzpicture}
